% Nota de custo_do_al. GERADA por flim_al/tabelas.py.
\footnotesize Todas as colunas são diferenças de F$\beta$ contra o \textbf{candidato sorteado}, medidas na validação, com a partição, as imagens, os marcadores base e o codificador inicial idênticos dentro de cada semente. \texttt{teto amostrado} é o melhor candidato entre os examinados, escolhido \textbf{pelo F$\beta$ da validação}: não é estratégia, e mede o que existe para capturar, não um teto absoluto. \texttt{o critério} é o que o Active Learning de fato escolhe, sem olhar resultado. A \textbf{lacuna} entre os dois é o custo de o critério errar o lugar do clique, e é a coluna com teste: t pareado por semente, com n = sementes, porque os candidatos de uma mesma semente compartilham codificador. \texttt{captura} é razão entre duas médias, fica instável com denominador pequeno e serve de ilustração, não de estatística. Linhas restritas à \textbf{base funcional}; a estratificação por regime é post-hoc e forçada pela aritmética, porque $\Delta$ a partir de F$\beta$=0 exato não pode ser negativo. As sementes compartilham pool e validação, então o IC vale para esta validação sob sorteio das imagens de treino e não generaliza para o dataset. O conjunto de teste não foi tocado por esta campanha.
\normalsize
